\documentclass[10pt,aspectratio=169]{beamer}
\input{../../_en_preambulo_beamer}
% Comandos y macros estadísticas compartidas para las presentaciones Beamer en inglés (Metropolis)

% Conjuntos numéricos básicos
\newcommand{\R}{\mathbb{R}}
\newcommand{\N}{\mathbb{N}}
\newcommand{\Q}{\mathbb{Q}}
\newcommand{\Z}{\mathbb{Z}}
\newcommand{\set}[1]{\left\{ #1 \right\}}
\newcommand{\abs}[1]{\left|#1\right|}
\newcommand{\norm}[1]{\left\| #1 \right\|}

% Comandos estadísticos y probabilísticos del curso
\newcommand{\Var}{\operatorname{Var}}
\newcommand{\cov}{\operatorname{Cov}}
\newcommand{\corr}{\rho}
\newcommand{\comb}[2]{\begin{pmatrix} #1 \\ #2 \end{pmatrix}}
\newcommand{\s}{\sigma}
\newcommand{\particion}{\mathfrak{P}}
\newcommand{\card}[1]{\#\left( #1 \right)}
\newcommand{\seq}[3]{\set{#1}_{#2}^{#3}}
\newcommand{\E}{\mathbb{E}}
\newcommand{\Prob}{\mathbb{P}}

% Entornos pedagógicos y cajas en inglés académico natural (soportando tanto nombres en inglés como en español)
\newenvironment{discussion}{%
  \begin{alertblock}{Discussion Question}%
}{%
  \end{alertblock}%
}
\newenvironment{discusion}{%
  \begin{alertblock}{Discussion Question}%
}{%
  \end{alertblock}%
}

\newenvironment{reflection}{%
  \begin{alertblock}{Classroom Reflection}%
}{%
  \end{alertblock}%
}
\newenvironment{reflexion}{%
  \begin{alertblock}{Classroom Reflection}%
}{%
  \end{alertblock}%
}

\newenvironment{paradox}{%
  \begin{alertblock}{Exploratory Paradox}%
}{%
  \end{alertblock}%
}
\newenvironment{paradoja}{%
  \begin{alertblock}{Exploratory Paradox}%
}{%
  \end{alertblock}%
}

\newenvironment{motivation}{%
  \begin{exampleblock}{Motivation: Data Science in Practice}%
}{%
  \end{exampleblock}%
}
\newenvironment{motivacion}{%
  \begin{exampleblock}{Motivation: Data Science in Practice}%
}{%
  \end{exampleblock}%
}

\newenvironment{quickchallenge}{%
  \begin{block}{Quick Team Challenge}%
}{%
  \end{block}%
}
\newenvironment{retoclase}{%
  \begin{block}{Quick Team Challenge}%
}{%
  \end{block}%
}


\title[Distributions of Functions]{Section 5.2: Probability Distributions of Functions of a Random Variable}
\subtitle{Preimages, level surfaces, and sums}
\author[J. Castillo Colmenares]{Juliho Castillo Colmenares}
\institute{Tecnológico de Monterrey}
\date{}

\begin{document}
\begin{frame}[plain]\vspace{-0.8cm}\titlepage\end{frame}

\begin{frame}{Roadmap --- Chapter 5}
  \footnotesize
  \begin{itemize}\itemsep=0.08cm
    \item \textbf{05.00} Inferential introduction
    \item \textbf{05.01} Transformation of variables
    \item \textbf{05.02} Functions of random variables \emph{(today)}
    \item \textbf{05.03--05.05} Classical sampling distributions
  \end{itemize}
  \begin{block}{Learning objective}
    Extend change of variable to non-monotone and multivariate functions, and
    recognize the Erlang distribution of a sum of exponentials.
  \end{block}
\end{frame}

\begin{frame}{Motivation: one function, one variable}
  \footnotesize
  When $Y=g(X)$, the density of $Y$ must gather every $X$ value producing the
  same result.
  \begin{alertblock}{Guiding question}
    How many preimages does each $y$ have, and how does the local scale change?
  \end{alertblock}
\end{frame}

\begin{frame}{One variable: general formula}
  \footnotesize
  For a differentiable function with regular preimages,
  \[
    f_Y(y)=\sum_{x:g(x)=y}f_X(x)\left|\frac{dx}{dy}\right|
    =\sum_{x:g(x)=y}\frac{f_X(x)}{|g'(x)|}.
  \]
  The sum runs over every solution of $g(x)=y$.
\end{frame}

\begin{frame}{Regularity condition}
  \footnotesize
  \begin{block}{Local hypothesis}
    $|g'(x)|\neq0$ at every preimage guarantees local invertibility there.
  \end{block}
  \begin{exampleblock}{Consequence}
    If the derivative vanishes, the critical point and resulting support need
    separate analysis.
  \end{exampleblock}
\end{frame}

\begin{frame}{Several variables: level surface}
  \footnotesize
  If $Y=g(X_1,\ldots,X_n)$ has a joint density,
  \[
    f_Y(y)=\int_{g(x_1,\ldots,x_n)=y}
    f_{X_1,\ldots,X_n}(x_1,\ldots,x_n)\frac{dS}{|\nabla g|}.
  \]
  The integral follows the level surface containing all configurations with the
  same value of $Y$.
\end{frame}

\begin{frame}{Independence and sums}
  \footnotesize
  For $Y=X_1+X_2$, the level surface is the line $x_1+x_2=y$. If the variables
  are independent, their joint density is the product of the marginal densities.
  \begin{alertblock}{Bridge}
    Integrating along that line produces the density of the sum.
  \end{alertblock}
\end{frame}

\begin{frame}{Reading of the section}
  \footnotesize
  \begin{itemize}\itemsep=0.12cm
    \item One variable: sum over preimages.
    \item Several variables: integrate over a level surface.
    \item The gradient $\nabla g$ measures local scale.
    \item Independence simplifies the joint density.
  \end{itemize}
\end{frame}

\begin{frame}{Derived application 1/4: find preimages}
  \footnotesize
  For each $y$, solve $g(x)=y$ before evaluating the formula.
  \begin{block}{Procedure}
    \begin{enumerate}\itemsep=0.1cm
      \item find all solutions $x$;
      \item check $g'(x)\neq0$;
      \item add the contributions.
    \end{enumerate}
  \end{block}
\end{frame}

\begin{frame}{Derived application 1/4: interpretation}
  \footnotesize
  \[
    \frac{f_X(x)}{|g'(x)|}
  \]
  is the contribution of one branch. The density of $Y$ is the sum of all
  branches reaching that value.
\end{frame}

\begin{frame}{Derived application 2/4: critical points}
  \footnotesize
  If $g'(x)=0$, the factor $1/|g'(x)|$ cannot be evaluated directly.
  \begin{alertblock}{Diagnostic}
    Critical points often mark support endpoints or density singularities and
    require a limit and total-mass check.
  \end{alertblock}
\end{frame}

\begin{frame}{Derived application 2/4: interpretation}
  \footnotesize
  The regularity condition is not merely formal: it prevents treating a map as
  invertible where it is not.
  \begin{exampleblock}{Reading rule}
    Identify the geometry of $g$ first; apply the formula second.
  \end{exampleblock}
\end{frame}

\begin{frame}{Solved application 3/4: exponential sum}
  \footnotesize
  Let $X_1,X_2\sim\operatorname{Exp}(1)$ be independent and let
  $Y=X_1+X_2$. For $y>0$, the level line has $0<x_1<y$ and $x_2=y-x_1$.
  \[
    f_{X_1,X_2}(x_1,x_2)=e^{-(x_1+x_2)}.
  \]
\end{frame}

\begin{frame}{Solved application 3/4: integration}
  \footnotesize
  \[
    f_Y(y)=\int_0^y e^{-(x_1+y-x_1)}\,dx_1
    =\int_0^y e^{-y}\,dx_1=ye^{-y},\quad y>0.
  \]
  \begin{alertblock}{Result}
    The sum has an Erlang density with shape $2$, equivalently
    $\operatorname{Gamma}(2,1)$.
  \end{alertblock}
\end{frame}

\begin{frame}{Solved application 4/4: generalization}
  \footnotesize
  The section concludes that the sum of $n$ iid exponential variables has
  \[
    \operatorname{Gamma}(n,1).
  \]
  The two-variable calculation exhibits how a level-surface integral produces a
  known family.
\end{frame}

\begin{frame}{Solved application 4/4: support}
  \footnotesize
  For the exponential sum, $X_1>0$ and $X_2>0$ imply $Y>0$.
  \begin{exampleblock}{Check}
    The density $ye^{-y}$ is interpreted for $y>0$ and integrates to one on
    that support.
  \end{exampleblock}
\end{frame}

\begin{frame}{Computational bridge 1/4: geometry}
  \footnotesize
  \begin{tabular}{ll}
    \hline
    Transformation & Set being integrated \\
    \hline
    $Y=g(X)$ & preimages $g^{-1}(y)$ \\
    $Y=g(X_1,X_2)$ & curve $g(x_1,x_2)=y$ \\
    $Y=g(X_1,\ldots,X_n)$ & level surface \\
    \hline
  \end{tabular}
\end{frame}

\begin{frame}{Computational bridge 2/4: joint density}
  \footnotesize
  For independent variables,
  \[
    f_{X_1,X_2}(x_1,x_2)=f_{X_1}(x_1)f_{X_2}(x_2).
  \]
  This factorization enters the integral over the level surface.
\end{frame}

\begin{frame}{Computational bridge 3/4: gradient}
  \footnotesize
  The term $|\nabla g|$ adjusts the surface element:
  \[
    dS/|\nabla g|.
  \]
  \begin{block}{Interpretation}
    The generalized Jacobian turns a geometric measure into probability mass.
  \end{block}
\end{frame}

\begin{frame}{Computational bridge 4/4: checklist}
  \footnotesize
  \begin{enumerate}\itemsep=0.12cm
    \item fix the support of $Y$;
    \item describe all preimages or level sets;
    \item write the joint density;
    \item integrate and check normalization.
  \end{enumerate}
\end{frame}

\begin{frame}{Synthesis of the section}
  \footnotesize
  \begin{itemize}\itemsep=0.12cm
    \item Preimages are added in one dimension.
    \item Level surfaces are integrated in several dimensions.
    \item The gradient controls scale.
    \item Sums of exponentials produce the Gamma/Erlang family.
  \end{itemize}
\end{frame}

\begin{frame}{Curricular transition}
  \footnotesize
  \begin{block}{Established}
    A function's distribution is obtained by gathering every sample point that
    produces the same transformed value.
  \end{block}
  \begin{alertblock}{Next section}
    \textbf{05.03 Chi-square distribution}: a classical sampling distribution
    built from squares of normal variables.
  \end{alertblock}
\end{frame}
\end{document}
